\begin{filecontents*}{\jobname.bib}
@book{BDB1907,
  author    = {Francis Brown and S. R. Driver and Charles A. Briggs},
  title     = {A Hebrew and English Lexicon of the Old Testament},
  year      = {1907},
  publisher = {Clarendon Press},
  location  = {Oxford}
}
@book{Gesenius1910,
  author    = {Wilhelm Gesenius and E. Kautzsch and A. E. Cowley},
  title     = {Gesenius' Hebrew Grammar},
  edition   = {2nd English},
  year      = {1910},
  publisher = {Clarendon Press},
  location  = {Oxford}
}
@book{Skinner1910,
  author    = {John Skinner},
  title     = {A Critical and Exegetical Commentary on Genesis},
  series    = {International Critical Commentary},
  year      = {1910},
  publisher = {T\&T Clark},
  location  = {Edinburgh}
}
@book{Speiser1964,
  author    = {E. A. Speiser},
  title     = {Genesis: Introduction, Translation, and Notes},
  series    = {Anchor Bible 1},
  year      = {1964},
  publisher = {Doubleday},
  location  = {Garden City, NY}
}
@book{Wenham1994,
  author    = {Gordon J. Wenham},
  title     = {Genesis 16--50},
  series    = {Word Biblical Commentary 2},
  year      = {1994},
  publisher = {Word Books},
  location  = {Dallas}
}
@book{Carr1996,
  author    = {David M. Carr},
  title     = {Reading the Fractures of Genesis: Historical and Literary Approaches},
  year      = {1996},
  publisher = {Westminster John Knox},
  location  = {Louisville}
}
@book{Nicholson1998,
  author    = {Ernest W. Nicholson},
  title     = {The Pentateuch in the Twentieth Century: The Legacy of Julius Wellhausen},
  year      = {1998},
  publisher = {Clarendon Press},
  location  = {Oxford}
}
@incollection{Roemer2006,
  author    = {Thomas R\"omer},
  title     = {The Elusive Yahwist: A Short History of Research},
  booktitle = {A Farewell to the Yahwist? The Composition of the Pentateuch in Recent European Interpretation},
  editor    = {Thomas B. Dozeman and Konrad Schmid},
  year      = {2006},
  publisher = {Society of Biblical Literature},
  location  = {Atlanta},
  pages     = {9--27}
}
@book{Friedman2003,
  author    = {Richard Elliott Friedman},
  title     = {The Bible with Sources Revealed},
  year      = {2003},
  publisher = {HarperSanFrancisco},
  location  = {San Francisco}
}
@book{AndersenFreedman1980,
  author    = {Francis I. Andersen and David Noel Freedman},
  title     = {Hosea: A New Translation with Introduction and Commentary},
  series    = {Anchor Bible 24},
  year      = {1980},
  publisher = {Doubleday},
  location  = {Garden City, NY}
}
@book{Aitken2007,
  author    = {James K. Aitken},
  title     = {The Semantics of Blessing and Cursing in Ancient Hebrew},
  series    = {Ancient Near Eastern Studies Supplement 23},
  year      = {2007},
  publisher = {Peeters},
  location  = {Leuven}
}
@book{Naveh1982,
  author    = {Joseph Naveh},
  title     = {Early History of the Alphabet: An Introduction to West Semitic Epigraphy and Palaeography},
  edition   = {2nd},
  year      = {1982},
  publisher = {Magnes Press},
  location  = {Jerusalem}
}
@article{Wimmer2023,
  author  = {Stefan Jakob Wimmer},
  title   = {Evidence of an Edomite Hebrew Scribal Cooperation from Tel Mal\d{h}ata: One Ostracon---Two Scripts---Three Scribes?},
  journal = {Atiqot},
  volume  = {110},
  year    = {2023},
  pages   = {33--43}
}
@incollection{Vanderhooft2014,
  author    = {David S. Vanderhooft},
  title     = {Iron Age Moabite, Hebrew, and Edomite Monumental Scripts},
  booktitle = {An Eye for Form: Epigraphic Essays in Honor of Frank Moore Cross},
  editor    = {Jo Ann Hackett and Walter E. Aufrecht},
  year      = {2014},
  publisher = {Eisenbrauns},
  location  = {Winona Lake, IN}
}
@book{Fleming2021,
  author    = {Daniel E. Fleming},
  title     = {Yahweh before Israel: Glimpses of History in a Divine Name},
  year      = {2021},
  publisher = {Cambridge University Press},
  location  = {Cambridge}
}
@article{MolenJacob,
  author  = {Steven L. Molen},
  title   = {The Identity of Jacob's Opponent: Wrestling with Ambiguity in Genesis 32},
  journal = {Dialogue: A Journal of Mormon Thought},
  note    = {online essay; accessed 2026}
}
@online{ORACCUgaritic,
  author  = {{Open Richly Annotated Cuneiform Corpus}},
  title   = {Ugaritic Glossary, entry \textit{\v{s}\char"02BEl}},
  url     = {https://oracc.museum.upenn.edu/aemw/ugarit/cbd/uga/SH.html},
  urldate = {2026-08-18}
}
@online{SAHDBRK,
  author  = {{Semantics of Ancient Hebrew Database}},
  title   = {Entry \textit{b-r-k}: to bless, praise},
  url     = {https://sahd-online.com/words/b-r-k_2/},
  urldate = {2026-08-18}
}
\end{filecontents*}

\documentclass[11pt]{article}

\usepackage{fontspec}
\usepackage{polyglossia}
\setdefaultlanguage{english}
\setotherlanguage{hebrew}
\setmainfont{FreeSerif}
\setsansfont{FreeSans}
\setmonofont{FreeMono}
\newfontfamily\hebrewfont[Script=Hebrew]{Noto Serif Hebrew}
\newcommand{\heb}[1]{\texthebrew{#1}}

\usepackage[a4paper,margin=27mm,headheight=15pt]{geometry}
\usepackage{microtype}
\usepackage{setspace}
\usepackage{booktabs}
\usepackage{tabularx}
\usepackage{longtable}
\usepackage{array}
\usepackage{enumitem}
\usepackage{csquotes}
\usepackage{fancyhdr}
\usepackage{titlesec}
\usepackage{footmisc}
\usepackage{xcolor}
\usepackage{hyperref}
\usepackage[backend=biber,style=authoryear,sorting=nyt,maxcitenames=2,maxbibnames=99]{biblatex}
\addbibresource{\jobname.bib}

\hypersetup{
  colorlinks=true,
  linkcolor=black,
  citecolor=black,
  urlcolor=black,
  pdftitle={The Wrong Brother: Esau's Defense and the Veiled Anti-Jacob Polemic of Genesis 25--33 and Hosea 12},
  pdfauthor={Esau Rothmann}
}

\setlength{\parindent}{1.25em}
\setlength{\parskip}{0.15em}
\setstretch{1.06}
\setcounter{secnumdepth}{3}
\setcounter{tocdepth}{2}
\emergencystretch=2em

\titleformat{\section}{\large\bfseries}{\thesection.}{0.55em}{}
\titleformat{\subsection}{\normalsize\bfseries}{\thesubsection.}{0.55em}{}
\titleformat{\subsubsection}{\normalsize\itshape}{\thesubsubsection.}{0.55em}{}

\pagestyle{fancy}
\fancyhf{}
\fancyhead[L]{\small Rothmann \textit{The Wrong Brother}}
\fancyhead[R]{\small \thepage}
\renewcommand{\headrulewidth}{0.3pt}

\newcommand{\versequote}[2]{%
  \begin{quote}\small
  \textit{#1}\\[-0.15em]
  #2
  \end{quote}
}
\newcommand{\translit}[1]{\textit{#1}}

\begin{document}

\begin{center}
{\LARGE\bfseries The Wrong Brother}\par
\vspace{0.35em}
{\large Esau's Defense and the Veiled Anti-Jacob Polemic of Genesis 25--33 and Hosea 12}\par
\vspace{1.2em}
{\large Esau Rothmann}\par
\vspace{0.15em}
{\itshape Olde Dominion University\footnote{Lit. ``Old Edominion University'' in the original.}}\par
\vspace{1.1em}
\end{center}

\begin{abstract}
Modern interpretation has tended to grant Jacob the moral prestige of the nation later named for him and to read Esau retrospectively through Israel--Edom hostility. This article argues that such reception reverses much of the literary evidence. In Genesis 25--33, Jacob bargains opportunistically for his exhausted brother's birthright, lies repeatedly to his father, impersonates Esau, takes the blessing intended for him, flees the consequences, and approaches the reunion convinced that the brother he wronged will behave as Jacob himself has behaved. Esau, when finally permitted to act, does the opposite: he runs to Jacob, embraces and kisses him, weeps with him, initially refuses compensation, calls him ``my brother,'' and offers assistance. The Jabbok scene is read here not as uncomplicated heroic praise of Jacob but as a veiled trial and humiliation of the national ancestor. Jacob must answer the identity question he once answered falsely, is renamed amid a conspicuous sequence of \heb{אמר} and \heb{שאל}, leaves injured, bows seven times before Esau, identifies Esau's face with the face of God, and finally commands the brother from whom he took the blessing, ``Please, take my blessing.'' Hosea 12 independently reads Jacob as an indictment of Israel and adds that the wrestler wept and supplicated. I argue that these features are best read as deliberate insider criticism: Israel's authors could praise their eponymous ancestor on the surface while allowing the narrative beneath that praise to vindicate Esau and expose Jacob. A final, more conjectural section considers graphic paronomasia between unpointed \heb{ישראל}, \heb{ישאל}, and \heb{עשו אל} within the mixed Hebrew--Aramaic--Edomite scribal ecology of the Iron Age. The thesis is unapologetically tendentious: later readers have defended the wrong brother.
\end{abstract}

\noindent\textbf{Keywords:} Jacob; Esau; Israel; Edom; Jabbok; Hosea 12; source criticism; \heb{שאל}; \heb{ברך}; palaeography; paronomasia; biblical satire

\section{Introduction: The National Ancestor Has Enjoyed Home-Field Advantage}

Jacob has enjoyed an interpretive advantage almost impossible to overstate: the community preserving the story eventually called itself by his new name. A national eponym does not enter reception history on equal terms with the neighboring ancestor whom that nation later opposed. ``Israel'' became self-designation; ``Edom'' became other. Once that political asymmetry is projected backward, Jacob's election tends to become moral exoneration and Esau's displacement tends to become proof that he deserved displacement.

The narrative itself is considerably less cooperative.

Jacob is introduced grasping his brother's heel. He later makes food conditional on the surrender of the birthright. He conspires with Rebekah to deceive a blind father, dresses himself in his brother's clothes, covers his skin to counterfeit his brother's body, and answers the question ``Who are you?'' with the categorical lie ``I am Esau your firstborn.'' He accepts the blessing intended for Esau and leaves. Esau's murderous intention after discovering the fraud is morally serious and should not be sanitized; yet it remains intention. When the brothers finally meet, the text narrates no assault, fraud, theft, curse, or act of vengeance by Esau. It narrates an embrace.

This article argues that this contrast is not accidental embarrassment within an otherwise Jacobite panegyric. It is the point. Genesis preserves a deliberately unstable national story in which the ``chosen'' brother repeatedly behaves worse than the ``rejected'' one, and in which his decisive transformation occurs only when he is compelled to stop pretending to be Esau, submit bodily before Esau, see God in Esau's face, and return to Esau what he had taken.

The argument is intentionally asymmetrical. Esau has had ample hostile reception; Jacob requires no additional counsel for the defense. The task here is to reconstruct the case the biblical authors themselves appear to have left for the other brother.

\section{J, E, Hosea, and the Possibility of Veiled Insider Criticism}

The labels J and E are used here heuristically. Pentateuchal source criticism has changed drastically since the classical Documentary Hypothesis, and contemporary scholars disagree over the dates, extent, independence, and even usefulness of the traditional sources \parencite{Carr1996,Nicholson1998,Roemer2006}. The argument therefore does not depend on assigning each sentence of Genesis 25--33 to a universally agreed hand. It requires only the more modest proposition that the Jacob cycle preserves multiple pre-priestly Israelite literary traditions, conventionally associated with J and E, and that these traditions were composed, transmitted, or edited by writers situated within Israelite or Judahite political worlds \parencite{Friedman2003,Carr1996}.

On the classical reconstruction, that situation is almost ideal for veiled national criticism. A Judahite writer and a northern Israelite writer could disagree on cult, geography, monarchy, and ancestral memory while sharing one dangerous rhetorical problem: Jacob/Israel was their own eponym. Direct denunciation of Jacob could be heard as denunciation of the nation. Superficial praise, however, could carry a second transcript beneath it. The ancestor can be ``chosen'' and still be ridiculous, dishonest, frightened, grasping, and in need of moral defeat.

This is not an exotic political technique. National literatures regularly authorize criticism by placing it in ancestral narrative, comedy, dream, proverb, prophetic sign, or deliberately ambiguous myth. The biblical corpus itself is unusually hospitable to internal dissent. Kings are mocked by prophets; royal ideology is preserved beside anti-monarchical speeches; Judah narrates its own disasters; Israel's heroes are rarely clean enough to function as straightforward propaganda.

The thesis advanced here is that the Jacob--Esau cycle belongs to this tradition of criticism from inside the house. Its surface legitimates Israel. Its subtext asks what sort of nation would choose Jacob as its ancestor and then dare to despise Esau.

Hosea is especially important because here the political critique is no longer hidden. Hosea invokes Jacob precisely while indicting Israel. The patriarch is not merely a revered founder; he is a precedent for the nation's habit of grasping, contending, bargaining, and deceiving. The older narrative ambiguity becomes prophetic accusation \parencite{AndersenFreedman1980}.

\section{The Case Against Jacob Begins Before the Jabbok}

\subsection{The birthright: hunger as an opportunity}

Genesis 25 famously closes the lentil episode by saying that Esau ``despised'' his birthright. That narrator's judgment is real. It is also not an acquittal of Jacob. The preceding action remains what it is: Jacob sees an exhausted brother asking for food and converts the request into a transaction for inheritance status. The text does not say Jacob generously feeds Esau and later negotiates; he makes eating contingent upon the oath.

A reader committed in advance to Jacob's superiority can make Esau's impulsiveness swallow the whole scene. A reader attending to conduct must notice that Esau's foolishness and Jacob's opportunism coexist. Election does not transform leverage over a hungry sibling into virtue.

\subsection{The blessing: identity theft in the literal sense}

Genesis 27 removes any remaining ambiguity about Jacob's methods. Rebekah plans the deception, but Jacob executes it. He enters Isaac's room wearing Esau's garments and animal skins designed to imitate Esau's body. Isaac's question is not obscure:

\versequote{Genesis 27:18--19}{``Who are you, my son?'' Jacob answers: ``I am Esau your firstborn.''}

Isaac remains suspicious and asks again whether he is really Esau. Jacob confirms the lie. He thereby receives the blessing intended for the absent brother. The scene is not metaphorical identity theft. It is identity theft performed by costume, tactile impersonation, false naming, and exploitation of a disabled parent.

Esau's later complaint is correspondingly precise: Jacob has taken his birthright and now ``taken my blessing.'' The verb of taking will matter when Genesis 33 reverses the transaction.

\subsection{Esau's guilt, and the striking limit the narrator places on it}

Esau is not spotless. Genesis reports that he despised the birthright, and after the stolen blessing he resolves to kill Jacob once Isaac dies. A defense of Esau that simply deletes those statements would deserve the same criticism this article directs at Jacob's later reception.

But the narrator draws a distinction modern summaries often erase: Esau's gravest wrongdoing remains contemplated and never narrated as action. Jacob's wrongdoing occupies scenes of accomplished deception. Esau threatens vengeance; Jacob has already defrauded him.

When Esau finally receives the opportunity to execute the threat, the text makes a point of narrating the opposite behavior. This is crucial. A character who once imagined murder approaches with four hundred men, finds the defrauding brother vulnerable, and then does not harm him. If the narrator wished to confirm Jacob's fear, the opportunity could scarcely have been better designed. Instead, Esau runs, embraces, kisses, and weeps.

The morally important Esau is therefore not the angry man imagining revenge in Genesis 27 but the man who has the power to take it in Genesis 33 and refuses.

\section{The Jabbok Is an Esau-Shaped Scene Before the ``Man'' Appears}

The anonymous wrestler is often approached as though he entered an empty symbolic landscape. He does not. Genesis 32 has spent the entire chapter constructing Esau as the event toward which Jacob is moving.

Jacob sends messengers to ``Esau my brother'' in Seir, the field of Edom. He learns that Esau is coming with four hundred men. He becomes terrified. He divides his camp. He prays for rescue ``from the hand of my brother, from the hand of Esau.'' He sends waves of gifts ahead of himself. Most revealingly, he rehearses the questions Esau is expected to ask.

The Hebrew introduces the anticipated speech with the root that will become conspicuous at the end of the wrestling scene:

\versequote{Genesis 32:18}{Esau will meet the servants and \heb{וּשְׁאֵלְךָ לֵאמֹר}---``ask you, saying ...''}

The imagined Esau asks questions of identity and ownership: Whose servant are you? Where are you going? Whose property is this? Thus, before the nocturnal opponent asks Jacob's name, Esau has already been scripted as an \emph{asker} of identity questions.

The chapter is equally saturated with ``face.'' Jacob hopes to appease Esau's face by an offering going before his own face, after which he will see Esau's face and perhaps Esau will lift his face. The Hebrew repetition of \heb{פנים} is conspicuous. The night will end with Jacob announcing that he has seen God ``face to face.'' The morning will end with Jacob telling Esau that seeing Esau's face is like seeing God's face.

The anonymous wrestler therefore enters a narrative corridor whose walls are already painted with Esau's questions and Esau's face. The man is unnamed; his literary environment is not.

\section{``Say Your Name'': The Identity Question as Judgment}

The decisive question at the Jabbok is not ``Can Jacob win a wrestling match?'' It is ``What is your name?''

That question repeats the moral crisis of Genesis 27. The last time Jacob was interrogated about his identity while a blessing was at stake, he said ``Esau.'' Now, again in bodily contact with another figure and again demanding a blessing, Jacob is asked:

\versequote{Genesis 32:28 [MT 32:28; English versification may differ]}{\heb{מַה־שְּׁמֶךָ}---``What is your name?''}

This time the answer is one word:

\versequote{}{\heb{יַעֲקֹב}---``Jacob.''}

The answer is routinely treated as mere narrative setup for the renaming. It is more plausibly a confession. Jacob has finally answered the old question truthfully.

The sequence can be stated without embellishment:

\begin{center}
\begin{tabularx}{0.88\textwidth}{@{}>{\bfseries}p{0.24\textwidth}X@{}}
\toprule
Genesis 27 & ``Who are you?'' --- ``I am Esau.'' --- Jacob receives Esau's blessing.\\
Genesis 32 & ``What is your name?'' --- ``Jacob.'' --- Jacob receives a new name.\\
\bottomrule
\end{tabularx}
\end{center}

Only after Jacob ceases pretending to be Esau may Jacob cease being Jacob.

This is why the opponent's refusal to disclose his own name is narratively apt. Jacob's problem is false identity; the opponent's function depends upon surplus identity. He is ``a man'' in Genesis, an angel in Hosea, encountered as God by Jacob, and immediately followed by Esau, whose face Jacob identifies with God's. The text forces Jacob into a single honest name while allowing his opponent to remain deliberately overdetermined.

\section{The \heb{אמר} Machine and the Sudden Appearance of \heb{שאל}}

The dialogue is constructed with unusual mechanical repetition. Verse numbering differs between traditions, so the sequence is more reliable than the printed numbers: after the bodily struggle, speech begins with repeated forms of \heb{אמר}, ``say.'' In common English versification the pattern runs through Genesis 32:26--29; in the Hebrew numbering it is shifted by one in some editions.

\begin{center}
\begin{tabularx}{\textwidth}{@{}p{0.16\textwidth}X@{}}
\toprule
\textbf{Unit} & \textbf{Speech architecture}\\
\midrule
Release / blessing & \heb{וַיֹּאמֶר} \dots{} \heb{וַיֹּאמֶר} --- ``and he said ... and he said''\\
Name question & \heb{וַיֹּאמֶר} \dots{} \heb{וַיֹּאמֶר} --- ``and he said ... and he said''\\
Renaming & \heb{וַיֹּאמֶר} \dots{} \heb{יֵאָמֵר} --- ``and he said ... it shall be said/called''\\
Climax & \heb{וַיִּשְׁאַל יַעֲקֹב וַיֹּאמֶר} \dots{} \heb{וַיֹּאמֶר} \dots{} \heb{תִּשְׁאַל}\\
\bottomrule
\end{tabularx}
\end{center}

The climax is lexically awkward in a useful way. The narrator does not merely write ``Jacob said.'' He first identifies the speech act:

\versequote{}{\heb{וַיִּשְׁאַל יַעֲקֹב וַיֹּאמֶר}---``And Jacob asked, and said ...''}

Jacob then says:

\versequote{}{\heb{הַגִּידָה־נָּא שְׁמֶךָ}---``Please, tell your name.''}

The answer repeats the marked root:

\versequote{}{\heb{לָמָּה זֶּה תִּשְׁאַל לִשְׁמִי}---``Why is it that you ask concerning my name?''}

After an almost drumlike succession of \emph{he said / he said / he said / he said}, the narrative suddenly insists on \emph{asking}. More remarkably, it does so immediately after the new name Israel has been pronounced.

The earlier Esau formula and the later Jacob formula now mirror one another:

\begin{center}
\begin{tabularx}{0.94\textwidth}{@{}p{0.45\textwidth}p{0.45\textwidth}@{}}
\toprule
\textbf{Expected Esau} & \textbf{Newly named Jacob}\\
\midrule
\heb{וּשְׁאֵלְךָ לֵאמֹר} & \heb{וַיִּשְׁאַל יַעֲקֹב וַיֹּאמֶר}\\
``he will ask you, saying'' & ``Jacob asked, and said''\\
\bottomrule
\end{tabularx}
\end{center}

The chapter first imagines Esau asking identity questions. Then the Esau-shaped stranger asks Jacob's identity. Then Jacob himself becomes the asker.

\section{\heb{ישאל}, \heb{ישראל}, and ``Please, God!''}

The lexical range of \heb{שאל} is broader than the English phrase ``ask a question.'' Biblical Hebrew uses it for asking, inquiring, asking for, requesting, demanding, consulting, and in some settings petitioning or praying for something \parencite[s.v. \heb{שאל}]{BDB1907}. Older grammatical descriptions use the contrast between Qal ``ask'' and an intensive Piel ``beg'' as an illustration of verbal intensification \parencite{Gesenius1910}. The Genesis forms in question are Qal, not Piel; a responsible argument must not quietly change their morphology. Yet context can make a Qal request desperate, and Genesis supplies an overt politeness/entreaty particle in the request itself: \heb{נָּא}, ``please.''

The Northwest Semitic comparison sharpens the visual pun. Ugaritic attests the cognate root \translit{\v{s}\char"02BEl}, including the prefix form \translit{y\v{s}\char"02BEal}, ``he asks/requests'' \parencite{ORACCUgaritic}. Hebrew gives the corresponding unpointed sequence:

\begin{center}
\begin{tabular}{@{}lll@{}}
\toprule
\textbf{Form} & \textbf{Consonantal string} & \textbf{Sense}\\
\midrule
\heb{ישאל} & y--\v{s}--\char"02BE--l & ``he asks / will ask''\\
\heb{ישראל} & y--\'s--r--\char"02BE--l & ``Israel''\\
\bottomrule
\end{tabular}
\end{center}

In the later square script, and more importantly in consonantal spelling without the Masoretic sin/shin dot, the written distinction is visually striking: \heb{ישאל} becomes \heb{ישראל} by the insertion of a single \heb{ר}. Historically the sibilants represented by later \heb{שׁ} and \heb{שׂ} were distinct, and the \emph{resh} is not dispensable. Thus ``Israel derives from \emph{yi\v{s}\char"02BEal}'' is not a defensible historical etymology.

As secondary paronomasia, however, the placement is almost ostentatious. Israel is named; the next marked action is ``Jacob asked''; his first request contains ``please''; the opponent answers with ``why do you ask?'' The text itself performs the pun it would need the reader to hear.

A deliberately joke-preserving English paraphrase is therefore stronger than a false etymology:

\begin{quote}\small
``Your name is Israel.''\\
And Jacob \emph{asks}: ``\emph{Please}, tell me your name.''\\
``Why do you \emph{ask} my name?''
\end{quote}

At its most compressed, the secondary echo is something like ``Ask El,'' ``Petition God,'' or the colloquial ``Please, God!'' The last is not a dictionary gloss; it is an English representation of the pragmatic cluster \heb{שאל} + \heb{נא} + divine referent. Its force lies in the comedy of the scene: the man whose official new name advertises that he has ``prevailed'' with God immediately behaves as a petitioner.

That is precisely the sort of double register expected if the surface celebrates Israel while the subtext humiliates Jacob.

\section{\heb{ברך}: Blessing, Knee, and the Joke the Body Finishes}

The dialogue is framed by \heb{ברך}. Jacob refuses release unless the opponent \emph{blesses} him:

\versequote{}{\heb{לֹא אֲשַׁלֵּחֲךָ כִּי אִם־בֵּרַכְתָּנִי}---``I will not release you unless you bless me.''}

The scene ends:

\versequote{}{\heb{וַיְבָרֶךְ אֹתוֹ שָׁם}---``And he blessed him there.''}

Both relevant verbal forms belong to the Piel. The standard lexical meaning here is ``bless,'' and no competent translation should replace the final clause with ``he forced him to kneel.'' Modern semantic work has further cautioned against assuming that the blessing root is historically identical with the root/noun complex associated with the knee and kneeling; Aitken and the Semantics of Ancient Hebrew Database treat the connection as unlikely \parencite{Aitken2007,SAHDBRK}.

Yet that historical caution does not eliminate a synchronic pun. Ancient speakers did not need a modern etymological tree to notice identical consonants. The language possessed \heb{ברך} vocabulary for blessing and BRK vocabulary associated with knees and kneeling. Older grammatical and lexicographical traditions often treated the Piel as an intensive or factitive formation and consequently made ``bless'' sound, at least to later readers, like an intensified bodily BRK \parencite{Gesenius1910,BDB1907}.

The exact formulation appropriate to the evidence is therefore:

\begin{quote}\small
Grammatically: ``and he \emph{blessed} him there.''\\
Punningly: ``and there he \emph{BRK-ed} him.''\\
Narratively: the next scene shows what being brought low looks like.
\end{quote}

This last point is decisive. Genesis does not require the disputed lexical connection to carry the argument. It literalizes the bodily submission immediately afterward.

\section{Seven Prostrations Before the Supposedly Inferior Brother}

At sunrise Jacob sees Esau approaching. Jacob goes ahead of his household and \heb{וַיִּשְׁתַּחוּ אַרְצָה שֶׁבַע פְּעָמִים}: he bows or prostrates himself to the ground seven times before reaching his brother.

The verb is not \heb{ברך}; it is the ordinary verb of prostration. That makes the literary effect stronger. The night scene ends with a potentially resonant ``BRK.'' The day scene supplies actual bodily abasement without lexical controversy.

The action also overturns one of the most politically charged lines in Isaac's stolen blessing. Jacob had obtained a blessing in which his mother's sons were told to bow before him. Yet when the two sons of Rebekah finally meet, it is Jacob who bows before Esau---not once, but seven times.

The stolen blessing predicts domination by Jacob. The enacted reconciliation stages submission by Jacob.

If this were straightforward ethnic propaganda, it is curiously bad propaganda.

\section{Esau's Face Is the Face of God}

The night scene closes with Jacob naming the place Peniel/Penuel because:

\versequote{}{\heb{כִּי־רָאִיתִי אֱלֹהִים פָּנִים אֶל־פָּנִים}---``For I have seen God face to face.''}

The following scene contains an unmistakable verbal echo. Jacob tells Esau:

\versequote{Genesis 33:10}{\heb{רָאִיתִי פָנֶיךָ כִּרְאֹת פְּנֵי אֱלֹהִים}---``I have seen your face as one sees the face of God.''}

This is not an inference imported from typology. Jacob himself supplies the comparison. The face he feared in Genesis 32 is the face he calls divine in Genesis 33.

Several modern readings have consequently asked whether the anonymous wrestler should be understood in relation to Esau rather than isolated from him; the ambiguity has been explored explicitly in literary interpretation \parencite{MolenJacob}. The most economical position is not that the opponent must be reduced to the historical person Esau. It is that Genesis deliberately overlays the figures. ``Man,'' angel, God, and Esauic double remain available simultaneously.

The theological force is uncomfortable for triumphalist readings of Jacob. Jacob does not finally discover God by defeating a stranger. He discovers that the brother he feared and defrauded bears the face of God.

\section{``Please, Take My Blessing'': The Transaction Is Reversed}

The argument reaches its clearest point in Genesis 33:11. Jacob first offers Esau a \heb{מנחה}, a present or offering. Esau initially refuses: he already has enough and tells Jacob to keep what is his. Jacob persists and changes the noun:

\versequote{}{\heb{קַח־נָא אֶת־בִּרְכָתִי}---``Please, take my blessing.''}

Translations sometimes smooth \heb{בִּרְכָתִי} into ``gift'' because material goods are in view. But the Hebrew says ``my blessing.'' That lexical choice is impossible to hear innocently after Genesis 27.

The reversal can be displayed almost algebraically:

\begin{center}
\begin{tabularx}{0.94\textwidth}{@{}p{0.44\textwidth}p{0.44\textwidth}@{}}
\toprule
\textbf{Genesis 27} & \textbf{Genesis 33}\\
\midrule
Jacob says ``I am Esau.'' & Jacob has confessed ``Jacob.''\\
Jacob takes Esau's blessing. & Jacob says, ``Take my blessing.''\\
Esau: ``he took my blessing.'' & Esau accepts Jacob's blessing.\\
Esau weeps. & The brothers weep together.\\
The blessing says brothers will bow to Jacob. & Jacob bows seven times to Esau.\\
\bottomrule
\end{tabularx}
\end{center}

The final verb is equally important: Esau takes it. What Jacob took under a false name is now taken back, with Jacob's repeated entreaty, under Jacob's true and transformed identity.

This is restitution in narrative form. No later moralizing speech is needed. The stolen blessing travels back across the boundary.

\section{Esau's Conduct in Genesis 33: The Text Refuses to Vindicate Jacob's Fear}

Jacob approaches the meeting as though Esau were Jacob. He assumes the brother will calculate, retaliate, exploit weakness, and seize what he can. Jacob divides people into defensive groups, sends wealth ahead, rehearses language, and prepares for catastrophe.

Esau does none of it.

The narrative gives him a rapid sequence of generous actions: he runs to meet Jacob, embraces him, falls on his neck, kisses him, and weeps. He asks about the family. He initially refuses the gift. He calls Jacob ``my brother.'' He offers to accompany him and then offers some of his people as protection or assistance.

The scene is remarkable because the author has designed every opportunity for Esau to become the villain later interpreters expect. Four hundred men are present. Jacob is frightened, encumbered with children and animals, and physically injured. The victim of fraud meets the perpetrator without legal institution, police, court, or neutral witness.

And Esau forgives him.

One may continue to cite Esau's old murderous intention; the text itself does. But to treat that intention as the essence of Esau while treating Jacob's accomplished frauds as incidental imperfections is not exegesis. It is national favoritism.

The reunion scene is written so that Jacob's moral imagination is exposed by contrast with Esau's actual conduct. The dangerous brother turns out to be the gracious one.

\section{The Graphic Innuendo: \heb{ישראל} as \heb{עשו אל}}

The most speculative part of the argument concerns not sound but sight.

In consonantal spelling, compare:

\begin{center}
\begin{tabular}{@{}lll@{}}
\toprule
\textbf{Reading} & \textbf{Consonantal analysis} & \textbf{Possible construal}\\
\midrule
\heb{ישראל} & y--\'s--r--\char"02BE--l & Israel\\
\heb{עשו אל} & \char"02BF--\'s--w + \char"02BE--l & Esau--El / Esau God\\
\bottomrule
\end{tabular}
\end{center}

Remove the word boundary from \heb{עשו אל} and the result is \heb{עשואל}. Three of the five consonantal positions are already identical. The remaining correspondences are \emph{yod/ayin} and \emph{resh/waw}.

This is not a historical etymology. The words do not sound sufficiently alike, the consonants are historically distinct, and no sound law derives Israel from ``Esau El.'' The proposal is graphic paronomasia: a learned visual joke that becomes available when related West Semitic scripts are placed in the same scribal horizon.

\subsection{What Edomite epigraphy contributes}

The Edomite epigraphic corpus is small, especially in comparison with Hebrew and Aramaic, and its monumental script is poorly represented. This limitation is fundamental \parencite{Vanderhooft2014}. What survives nevertheless shows precisely the sort of borderland scribal ecology relevant to the proposal. Edomite hands participate in a West Semitic continuum and show substantial Aramaic influence; the eastern Negev in particular preserves Hebrew and Edomite material in close contact. Naveh's classic palaeographic synthesis describes the gradual Aramaization of Edomite writing \parencite{Naveh1982}.

Recent work on an ostracon from Tel Mal\d{h}ata goes further, proposing Hebrew and Edomite scribal cooperation in a single epigraphic context \parencite{Wimmer2023}. Whatever one decides about individual readings, the broader historical point is secure: Judahite, Edomite, and Aramaic writing were not hermetically sealed visual systems.

The stronger of the two proposed graphic substitutions is \emph{resh/waw}. Actual West Semitic epigraphy can produce signs for which a rounded or abbreviated \emph{waw} and a \emph{resh}-like form become competing readings, especially in cursive, damaged, or regionally influenced hands. The weaker substitution is \emph{yod/ayin}. In formal monumental alphabets the signs are ordinarily distinguishable. Aramaizing cursive forms can simplify and open \emph{ayin}, while \emph{yod} can vary drastically, but no presently identified Iron Age Edomite hand has been demonstrated in which the entire pair \heb{ישראל}/\heb{עשואל} becomes a natural scribal confusion.

This point should be conceded, not concealed. The visual hypothesis is not evidence for a lost manuscript reading. It is a proposal about deliberate cross-script resemblance. A pun may exploit similarities that would be inadequate to explain a copying error.

\subsection{Why the proposed joke is nevertheless suspicious}

A graphic resemblance matters only if the narrative supplies the semantic key. Here it does.

The hypothetical resegmentation produces ``Esau--El.'' The very next scene makes Jacob say that seeing Esau's face is like seeing the face of God. The proposed letters and the prose therefore point in the same direction.

The joke, if intentional, would be unusually economical:

\begin{quote}\small
Officially: Israel, the one who strives with God.\\
Lexically nearby: \heb{ישאל}, the one who asks.\\
Graphically lurking: \heb{עשו אל}, Esau--El.\\
Narratively explicit: ``seeing your face is like seeing the face of God.''
\end{quote}

No single item proves the others. The accumulation is the argument.

\section{YHWH Comes from the Direction of Esau}

The Esau--divine overlap also participates in a larger biblical geography. Several archaizing Yahwistic poems imagine YHWH arriving from the southern regions associated with Seir, Edom, Teman, Paran, and Sinai. Judges 5 depicts YHWH going out from Seir and marching from the field of Edom. Deuteronomy 33 has YHWH come from Sinai, dawn from Seir, and shine from Paran. Habakkuk 3 has God come from Teman and the Holy One from Paran.

These poems are central evidence in modern discussions of a southern background for early Yahwism, though the historical reconstruction remains contested and should not be reduced to the formula ``YHWH was simply an Edomite god'' \parencite{Fleming2021}. The poetic datum is firmer than any single origin theory: Israel's oldest divine geography repeatedly places YHWH's dramatic arrival to Israel's south and southeast, in the same broad landscape associated with Edom and Seir.

Genesis 32 is therefore suggestive in a way later doctrinal boundaries can obscure. Esau is explicitly coming from Seir, the field of Edom. Before he becomes visible, a divine or angelic ``man'' arrives in the night. At sunrise Jacob sees Esau and says that Esau's face is like God's face.

The sequence is not a proof of cultic identity. It is literary alignment:

\begin{quote}\small
YHWH comes from Seir/Edom in ancient poetry.\\
Esau comes from Seir/Edom in Genesis 32.\\
A divine opponent appears immediately before Esau.\\
Jacob then recognizes God's face in Esau's face.
\end{quote}

For a text already willing to let divine identity shimmer between man, angel, and God, the Edomite brother is not an arbitrary fourth term.

\section{Hosea 12: The Ancient Author Who Gives the Game Away}

Hosea 12 is the strongest control because it shows a biblical author using Jacob not to flatter Israel but to prosecute it.

Hosea recalls Jacob grasping his brother in the womb and then says that in his strength he strove with God. The next clause reads:

\versequote{Hosea 12:5 in common English numbering}{\heb{וַיָּשַׂר אֶל־מַלְאָךְ וַיֻּכָל}---``He strove with an angel and prevailed.''}

At the level of consonantal writing, the opening words are remarkable:

\begin{center}
\heb{וישר אל} \quad $\longleftrightarrow$ \quad \heb{וישראל}
\end{center}

Remove the word boundary and ``and he strove with/to El'' becomes the exact consonantal sequence ``and Israel.'' The point is not that Hosea proposes a different historical etymology. The point is methodological. An ancient biblical author discussing this very story demonstrably treats the Israel sequence as a string that can be segmented into meaningful verbal material.

Hosea then adds the line Genesis does not state explicitly:

\versequote{}{\heb{בָּכָה וַיִּתְחַנֶּן־לוֹ}---``he wept and supplicated him.''}

The subject and object have generated interpretive debate, and some traditional readers reverse the direction of the supplication. For the present argument either possibility is damaging to a simplistic triumphal reading. Hosea remembers the ``prevailing'' wrestling episode as a scene involving tears and begging \parencite{AndersenFreedman1980}.

The relation to Genesis 32 is striking:

\begin{center}
\begin{tabularx}{0.93\textwidth}{@{}p{0.44\textwidth}p{0.44\textwidth}@{}}
\toprule
\textbf{Genesis 32} & \textbf{Hosea 12}\\
\midrule
Israel named through striving & Israel sequence resegmented through ``he strove + El''\\
\heb{וישאל} --- ``and he asked'' & ``he wept''\\
\heb{נא} --- ``please'' & ``he supplicated/begged''\\
\heb{תשאל} --- ``why do you ask?'' & plea for favor\\
Jacob ``prevails'' but leaves injured & Jacob ``prevails'' while weeping/supplicating\\
\bottomrule
\end{tabularx}
\end{center}

Hosea thus confirms two components of the proposed subtext independently: the Israel-name is available for graphic/verbal play, and the wrestling tradition is compatible with supplication rather than macho conquest.

More importantly, Hosea's purpose is anti-triumphal. He invokes Jacob while accusing contemporary Israel of deceit, false balances, misplaced confidence, and the repetition of ancestral crookedness. The national ancestor functions as a mirror in which the nation sees its vice.

If Genesis hid the criticism under the honor of election, Hosea removes the polite covering.

\section{Jacob Is Made to Become Esau Before He Is Allowed to Become Israel}

Hosea's detail of weeping and supplication creates another symmetry. In Genesis 27, after discovering the theft, Esau cries out bitterly, pleads ``Bless me---me too, my father,'' and weeps. In the Jabbok tradition as Hosea remembers it, Jacob is the one who weeps and pleads in a struggle over blessing.

The positions have reversed:

\begin{center}
\begin{tabularx}{0.95\textwidth}{@{}p{0.45\textwidth}p{0.45\textwidth}@{}}
\toprule
\textbf{Esau after the theft} & \textbf{Jacob at the Jabbok}\\
\midrule
Deprived of blessing & Demands blessing\\
Cries out & Struggles and cries in Hosea\\
Weeps & Weeps in Hosea\\
Pleads to be blessed & Supplicates / asks ``please''\\
Identity stolen by Jacob & Jacob must state his own identity\\
\bottomrule
\end{tabularx}
\end{center}

Jacob once became Esau by costume in order to steal Esau's blessing. The Jabbok makes him become Esau morally: the desperate petitioner who cannot command a blessing and must ask for one.

Only after experiencing the victim's position does Jacob become Israel.

This is a much more coherent transformation than the conventional image of a hero rewarded for beating God in a night fight. The limp is not a trophy. The new name is not merely a medal. The scene subjects the deceiver to a symbolic reversal of his crime.

\section{A Political Reading: Praise as Cover for Criticism of Israel}

What would such a story do for an Israelite or Judahite author?

It would permit national criticism without surrendering national identity. Jacob remains the ancestor. Israel remains the name. The blessing remains. Election is not revoked. Yet the very narrative that legitimates the nation also prevents the nation from imagining that election means moral superiority.

That tension makes excellent sense whether one imagines J and E as products of neighboring kingdoms, as traditions emerging from the memory of a united monarchy, or as later literary complexes preserving northern and southern perspectives. In every version, Jacob is politically dangerous material because Jacob is Israel. To criticize him is to criticize oneself.

The story therefore uses what may be called \emph{protective praise}. Jacob receives the surface markers required of a national ancestor: promises, blessing, divine encounter, a new name, and the statement that he ``prevailed.'' Under those honors, however, the author places a running indictment:

\begin{enumerate}[leftmargin=2em,itemsep=0.2em]
\item Jacob's cleverness repeatedly takes the form of exploitation or deception.
\item His greatest blessing is acquired under the stolen identity ``Esau.''
\item He fears Esau partly because he projects his own transactional morality onto him.
\item The divine encounter forces him to answer his own name truthfully.
\item The triumphal name is followed immediately by asking, ``please,'' and more asking.
\item The ``prevailer'' leaves physically damaged.
\item He prostrates himself seven times before the brother supposedly subordinated by the blessing.
\item He identifies that brother's face with God's face.
\item He explicitly returns ``my blessing'' to the man from whom he took a blessing.
\item Hosea uses Jacob as evidence against Israel and remembers the struggle as involving weeping and supplication.
\end{enumerate}

This is not what clumsy ethnic propaganda looks like. It is what a sophisticated national satire looks like.

The official message is safe: Israel prevailed with God.

The dangerous message is legible to anyone paying attention: Israel survives by finally learning to stop behaving like Jacob.

\section{The Esau Defense: What the Story Actually Narrates}

At this point the imbalance in reception can be stated plainly.

Jacob deceives his father. Esau does not.

Jacob impersonates his brother. Esau does not.

Jacob takes the blessing intended for Esau. Esau does not retaliate by taking Jacob's blessing; Jacob begs him to accept it.

Jacob anticipates violence. Esau embraces him.

Jacob addresses Esau as ``my lord.'' Esau answers ``my brother.''

Jacob expects to purchase peace. Esau initially refuses payment.

Jacob bows repeatedly. Esau does not demand the bows.

Jacob says Esau's face is like God's. Esau does not say the same of himself.

To call Esau ``always nice'' would be philologically imprecise if it erased the murderous intention in Genesis 27. To say that the narrator carefully writes Esau's \emph{actual conduct toward Jacob} as consistently more gracious than Jacob expects is simply to report the sequence.

The rehabilitation of Esau therefore requires surprisingly little special pleading. One need only stop converting Jacob's election into moral approval.

Indeed, a writer seeking covertly to defend Edom could hardly design the reunion more effectively. The Edomite ancestor arrives with overwhelming force and refuses vengeance. Israel's ancestor arrives with gifts and fear, bows to the earth, calls Edom ``lord,'' recognizes God in Edom's face, and returns the blessing.

The national hierarchy survives formally. The moral hierarchy has been inverted.

\section{Objections}

\subsection{``But Genesis says Esau despised the birthright''}

It does. Esau's impulsiveness is part of the characterization. The verse does not make Jacob's exploitation disappear. Two moral judgments can occupy one scene.

\subsection{``But Esau planned to kill Jacob''}

He did. That is Esau's strongest negative evidence. The narrative's later refusal to enact that intention is therefore all the more meaningful. Forgiveness is only narratively interesting where vengeance was possible.

\subsection{``Israel is explicitly explained from \heb{שׂרה}, not \heb{שאל} or Esau''}

Correct. The argument concerns secondary wordplay, not replacement etymology. Biblical naming scenes frequently sustain more than one phonetic or semantic association. Hosea's \heb{וישר אל}/\heb{וישראל} demonstrates that the Israel string itself invited resegmentation in ancient biblical writing.

\subsection{``\heb{שאל} does not literally mean `Please, God!' ''}

Correct. ``Please, God!'' is a literary paraphrase of a cluster: ask/request + \heb{נא} ``please'' + a divine context. The Qal forms in Genesis should not be relabeled Piel. The claim is pragmatic and paronomastic.

\subsection{``The Piel of \heb{ברך} means bless''}

Correct again. The paper translates it ``bless.'' The kneeling proposal is an undertone, strengthened by the bodily narrative that follows, not an alternative morphological parse. Modern lexical scholarship may separate the historical roots entirely \parencite{Aitken2007,SAHDBRK}.

\subsection{``The palaeographic Esau--El pun is too speculative''}

This is the best objection. The complete visual equivalence has not been demonstrated in one securely dated hand. The \emph{resh/waw} leg is stronger than the \emph{yod/ayin} leg. The thesis can stand without the pun; the pun becomes interesting because the surrounding narrative independently supplies Esau--God association.

\subsection{``J and E may not exist in the classical form presupposed here''}

The paper uses the labels as shorthand for pre-priestly or non-priestly ancestral traditions and for the longstanding observation that Genesis is composite. If future models rename or redistribute those materials, the literary argument remains: multiple Israelite biblical voices, culminating unmistakably in Hosea, preserve anti-triumphal readings of Jacob.

\section{Conclusion: Later Readers Defended the Wrong Brother}

The story of Jacob and Esau has generally been read from the standpoint of the winner's archive. Jacob becomes Israel; Israel preserves the book; Esau becomes Edom; later hostility supplies the moral coloring. Once the political outcome is allowed to interpret the ancestral story, Jacob's selection becomes virtue and Esau's displacement becomes guilt.

The narrative itself resists that simplification at nearly every decisive point.

Jacob acquires advantage through bargaining and disguise. He says ``I am Esau'' to obtain Esau's blessing. He flees. When he returns, he assumes Esau will behave violently and transactionally. Before Esau arrives, an anonymous figure intercepts him in an Esau-saturated chapter. A blessing is again at stake. An identity question is again asked.

This time Jacob tells the truth.

``Jacob.''

Only then is he renamed.

The new name is officially glorious: Israel, because he has striven with God and humans and prevailed. Yet the glory is immediately destabilized. The speech machinery of \heb{אמר} suddenly becomes \heb{שאל}. Jacob asks. He says ``please.'' The opponent answers, ``Why do you ask?'' The one who has supposedly prevailed leaves limping. The scene ends with a Piel \heb{ברך}: unmistakably ``he blessed him,'' but resonant within a story that will soon put Jacob on the ground.

Morning supplies the commentary. Jacob bows seven times before Esau. Esau runs to him, embraces him, kisses him, and weeps. Jacob sees Esau's face and says he has seen the face of God. Then Jacob says the sentence that resolves Genesis 27:

\begin{quote}\small
\heb{קַח־נָא אֶת־בִּרְכָתִי}.\\
``Please, take my blessing.''
\end{quote}

And Esau takes it.

Hosea later makes explicit what Genesis could veil. The Israel-name is visually exposed in \heb{וישר אל}; the patriarch becomes an indictment of Israel; and the victorious wrestling tradition is remembered with tears and supplication. The nation called Israel is told, in effect, that its ancestor's defining struggle was not the moment he proved stronger than God but the moment he was forced to plead, confess, limp, and change.

The speculative graphic joke \heb{ישראל} $\sim$ \heb{עשו אל} belongs at the edge of the case, not its center. Yet it is precisely the kind of edge on which biblical innuendo often lives. Its palaeographic support is incomplete; its narrative support is almost indecently convenient. ``Esau--El'' is what the letters might be made to suggest. ``Your face is like the face of God'' is what the text actually says.

The political implication is sharper than the usual reconciliation sermon. The writers responsible for Israel's ancestral literature may have protected criticism of their own polity beneath the superficial honor due its eponym. Jacob could be elected and still be the butt of the story. Esau could be displaced and still be the morally better brother. Edom could be the foreign other while also supplying the face in which Israel finally learns to recognize God.

Such a reading does not abolish Israel's blessing. It tells us what the blessing costs.

Jacob may keep the name Israel only after he stops stealing the name Esau.

Israel may keep the blessing only after Jacob offers the blessing back.

And Israel may claim to have seen God only after Jacob can look at Esau and admit where God's face was all along.

\section*{Competing Interests}

The author declares no competing interests, apart from a longstanding familial concern regarding the disposition of one paternal blessing and an understandable sensitivity to later national caricatures of Edom.

\section*{Acknowledgments}

The author thanks the anonymous reviewers who, unlike certain relatives, returned what did not belong to them.

\printbibliography[title={References}]

\end{document}
